% Part: computability
% Chapter: computability-theory
% Section: ce-sets

\documentclass[../../../include/open-logic-section]{subfiles}

\begin{document}

\olfileid{cmp}{thy}{ces}
\olsection{المجموعات القابلة للتعداد حسابيًا}

\begin{defn}
تكون المجموعة \emph{قابلة للتعداد حسابيًا} إذا كانت خالية أو كانت مدى
دالة قابلة للحساب.
\end{defn}

\begin{history}
تسمى المجموعات القابلة للتعداد حسابيًا أيضًا \emph{قابلة للتعداد
عوديًا}. وهذه هي التسمية الأصلية، ويشيع اليوم استعمال التسميتين، وكذلك
الاختصارين ``c.e.'' و``r.e.''.
\end{history}

\begin{explain}
ينبغي أن تفكر فيما يعنيه التعريف ولماذا تلائم التسمية. والفكرة هي أنه
إذا كانت $S$ مدى الدالة القابلة للحساب~$f$، فإن
\[
S = \{ f(0), f(1), f(2), \dots \},
\]
ومن ثم يمكن النظر إلى $f$ على أنها «تعدد» عناصر~$S$. ولاحظ أن التعريف
لا يشترط أن تكون~$f$ دالة متزايدة، أي لا يلزم أن يكون التعداد مرتبًا
ترتيبًا متزايدًا. بل لا يلزم حتى أن تكون $f$ أحادية، أي يسمح بالتكرار
في تعداد~$S$ المعطى بـ$f(0)$ و$f(1)$ و$f(2)$ و…. فمثلًا تعدد الدالة
الثابتة $f(x)=0$ المجموعة $\{0\}$.
\end{explain}

كل مجموعة قابلة للحساب قابلة للتعداد حسابيًا. ولرؤية ذلك، افترض أن
$S$ قابلة للحساب. فإذا كانت $S$ خالية، كانت قابلة للتعداد حسابيًا
بحسب التعريف. وإلا فليكن $a$ أي عنصر من $S$. عرف $f$ بـ
\[
f(x) =
\begin{cases}
x & \text{إذا كانت $\Char{S}(x) = 1$} \\
a & \text{في غير ذلك.}
\end{cases}
\]
عندئذ تكون $f$ دالة قابلة للحساب، وتكون $S$ مدى~$f$.

\end{document}
